\documentclass[10pt]{article}

\usepackage[utf8]{inputenc}

\usepackage[T1]{fontenc}

\usepackage{amsmath}

\usepackage{amsfonts}

\usepackage{amssymb}

\usepackage[version=4]{mhchem}

\usepackage{stmaryrd}

\usepackage{mathrsfs}



\begin{document}

Блиновой. Роль вихря двумерного течения здесь играет т. н. потенциальный вихрь - скалярное произведение вихря абсолютной скорости на градиент энтропии. Уравнение для него в квазигеострофич. приближении получается в виде



\[

\frac{\partial \mathscr{L} \psi}{\partial t}+J(\psi, \mathscr{L} \psi)+\beta \frac{\partial \psi}{\partial x}=F ; \mathscr{L}=\Delta+\frac{\partial}{\partial z} \frac{H^{2}}{L_{R}^{2}} \frac{\partial}{\partial z},

\]



где $\psi, \Delta, J$ - горизонтальные функция тока, лапласиан и якобиан, $z$ - вертикальная координата, $H$ и $L_{R}=H N / f$ - толщина слоя и «радиус деформации» Россби ( $N$ - частота Вайссала - Брента, $f$ - параметр Кориолиса), $x$ - дуга круга широты, $\beta$ - производная от $f$ по дуге меридиана, $F$ - неадиабатические факторы. В масштабах $L \ll L_{R}$ получается обычное уравнение двумерной турбулентности. Волновое число $k_{\beta}=(\beta / 2 U)^{1 / 2}$ (где $U$ - типичная скорость) разделяет вихри $\left(k>k_{\beta}\right)$ и волны Россби - Блиновой $\left(k<k_{\beta}\right)$. Малые начальные вихри имеют тенденцию со временем расти, выравниваться іо вертикали («баротропизация»), смсщаться на запад и вытягиваться вдоль кругов широты.



Важными обобщениями турбулентности в несжимаемой жидкости являются турбулентности в страфицированной жидкости с архимедовыми силами и гидромагнитная турбулентность.



Лит.: [1] М о н и н А. С., Я г ло м А. М., Статистическая гидромеханика, ч 1-2, м., 1965-67, L2] Г л е д з е p E. Б., М о н и н А. С., “Успехи матем. наую», 1974, т. 29, в. 3, с. 11159; [3] М о н и н А. С., "Успехи физич. наук", 1978, т. 125 , в. 1, c. 97-122, [4] и и "У спех м о в Р. В., Океанская турбулентность, Л., 1981; [6] О з м и д о в Р. В., Океанская турбулентность, Л., 1981; [6] В и ши и к М. И., Ф Ф р с и к о в А. В., Математические зада-

чи статистическй гидромеханики, М., 1980.



ТУРНИР - ориентированный граø без петель, каждая пгара вершин к-рого соединена дугой точно в одном направлении. Т. с $n$ верпинами может служить описанием исхода состязания $n$ игроков, правилами к-рого запрещен ничейный исход. Понятие Т. исіголзуется для упорядочения $n$ объектов методом попарных сравнений. В связи с этим оно находит свои приложения в биологии, социологии и т. п.



Т. наз. т р а н 3 и т и в ны м, если можно так занумеровать его вершины числами $1,2, \ldots, n$, что из вершины $v_{i}$ идет дуга в вершину $v_{j}$ тогда и только тогда, когда $i>j$. В транзитивном Т. отсутствуют контуры. Т. наз. с и л ь н ы м, если для любой упорядоченной пары его вершин $v_{i}, v_{j}$ существует ориентированный путь из $v_{i}$ в $v_{j}$. Множество дуг в Т. наз. с о гл а с о в а н ны м, если в подграфе, образованном этими дугами и инцидентными им вершинами, отсутствуют контуры. Максимальная мощность множества согласованных дуг является мерой согласованности при определении «победителя» Т. Всякий Т. содержит подмножество согласованных дуг мощности не меньшей, чем $\left(n^{2 / 4}\right)(1+o(1))$. Другой мерой согласованности является отношение числа транзитивных $\boldsymbol{\kappa}$-верпинных подтурниров турнира с $n$ вершинами к числу сильных $k$-вершинных подтурниров. Максимальное число сильных $k$-вершинных подтурниров турнира c $n$ вершинами равно



\[

\begin{gathered}

C_{n}^{k}-k C_{\frac{n-1}{2}}^{2}, \text { еслии } n \text { нечетно, } \\

C_{n}^{k}-\frac{k}{2}\left(C_{\frac{n}{2}}^{k+1}-C_{\frac{n-2}{2}}^{k-1}\right), \text { если } n \text { четно. }

\end{gathered}

\]



Т. является сильным тогда и только тогда, когда он имеет остовный контур. Каждый сильный $\mathrm{T}$. с $n$ верциинами имеет контур длины $k$ для $k=3,4, \ldots, n$. Всякий Т. имеет остовный путь. Набор всех полустегеней исхода $d_{i}$ турнира с $n$ вершипами удовлетворяет равенству



\[

\sum_{i=1}^{n} d_{i}^{2}=\sum_{i=1}^{n}\left(n-d_{i}\right)^{2}

\]



Пусть набор целых чисел $\left(d_{1}, \ldots, d_{n}\right)$ удовлетворяет условию $0 \leqslant d_{1} \leqslant \ldots \leqslant d_{n} \leqslant n-1$. Т. с полустепенями исхода $d_{1}, \ldots, d_{n}$ существует тогда и только тогда, когда для любого $k=1, \ldots, n-1$ выполняется неравенство



\[

\sum_{i=1}^{k} d_{i} \geqslant \frac{k(k-1)}{2}

\]



а для $k=n$ справсдливо равенство. При этом Т. является сильным тогда и только тогда, когда



\[

\sum_{i=1}^{k} d_{i}>\frac{k(k-1)}{2} \text { для } k<n .

\]



Если $T_{1}$ и $T_{2}$-два подтурнира турнира $T$ и для каждой пары вершин $v^{\prime}$ из $T_{1}, v^{\prime \prime}$ из $T_{2}$ существует дуга $\left(v^{\prime}, v^{\prime \prime}\right)$, то пишут $T_{1} \rightarrow T_{2}$. Пусть множество вершин Т. разбито на непересекающиеся подмножества $T_{1}, \ldots$, $T_{m}$. И пусть либо $T_{i} \rightarrow T_{j}$, либо $T_{j} \rightarrow T_{i}$ для $1 \leqslant i \leqslant j \leqslant$ $\leqslant m$. Тогда разбиение определяет о т но ші е в и е ә к в и в а л е н т н о с т и на вершинах Т. T. Турнир наз. п р о с т ы м, если на его вершинах нельзя определить нетривиальное отношение әквивалентности. Всякий Т. с $n$ вершинами являөтся подтурниром нек-рого простого T. с $n+2$ вершинами. Турнир $T$ с $n$ вершинами является подтурниром нек-рого простого Т. с $n+1$ верпинами тогда и только тогда, когда $T$ не является ни контуром с тремя вершинами, ни нетривиальным транзитивным T. с нечетным числом верпин. Число попарно неизоморфных Т. с $n$ вершинами асимптотически равно $\frac{1}{n !} 2^{C_{n}^{2}}$. Число различных $\mathrm{T}$. с $n$ нумерованными вершинами равно $2^{C_{n}^{2}}$. Производяпцие функции $t(x)$ и $s(x)$ для Т. и сильных Т. соответственно связаны соотноптоием:



\[

s(x)=\frac{t(x)}{1+t(x)} .

\]



Зсякий Т. с $n$ вершинами, $n \geqslant 5$, не являющийся сильным, однозначно восстанавливается по совокупности своих вершинных подтурниров.



Лит.: [1] X а р а р и Ф., Теория графов, пер. с англ., М., 1973, [2] M o o n J. W., Topics on tournaments, N. Y., [1968].



ТУЭ МЕТОД - метод в теории диофантовых приближений, созданный А. Туэ [1] в связи с проблемой приближсиия алгебраич. чисел рациональными числами: найти велитину $v=v(n)$, при к-рой для каждого алгебраич. числа $\alpha$ степени $n$ неравенство



\[

\left|\alpha-\frac{p}{q}\right|<\frac{1}{q^{v+\varepsilon}}

\]



имеет конетное число решений в целых рациональных числах $p$ и $q, q>0$, при любом $\varepsilon>0$, и бесконечное число решений при любом $\varepsilon<\Theta$.



А. Туә показал, что $v \leqslant \frac{n}{2}+1$. Т. м. основан на свойствах специального многочлена $f(x, y)$ от двух переменных $x, y$ с целыми коэффициентами и предположении существования двух решений неравенства (1) при $v \leqslant \frac{n}{2}+1$ с достаточно больпими значениями $q$. Т е oр е м а Т у ә имеет много важных приложений в теории чисел. В тастности; из нее следует, что диофантово уравнение



\[

F(x, y)=m,

\]



где $F(x, y)$ - неприводимая форла от переменных $x$ и $y$ с целыми коәффициентами степени $n \geqslant 3$, а $m$ -





\end{document}